\documentclass[10pt]{extarticle}
\usepackage[a4paper,margin=1.5cm]{geometry}
\usepackage[T1]{fontenc}
\usepackage{amsmath,amssymb}
\usepackage{graphicx}
\usepackage[export]{adjustbox}
\usepackage{multicol}
\usepackage{enumitem}
\usepackage{booktabs,array,tabularx}
\setlength{\columnsep}{0.7cm}
\setlength{\parindent}{0pt}
\setlength{\parskip}{3pt}
\pagestyle{empty}
\makeatletter
% figure registry (figdefs.tex is generated by the extract step)
\newcommand{\qfigdef}[4]{\@namedef{qf@file@#1}{#2}\@namedef{qf@w@#1}{#3}\@namedef{qf@wide@#1}{#4}}
\newcommand{\qf@check}[1]{\@ifundefined{qf@file@#1}{\PackageError{score2}{Unknown figure id #1}{}}{}}
\newcommand{\qf@img}[4][]{\includegraphics[width=#2\dimexpr\@nameuse{qf@w@#4}pt\relax,max width=#3,max height=0.45\textheight,#1]{\@nameuse{qf@file@#4}}}
% Full-width questions: finalize wraps every question that contains a wide figure in \qfullbegin...\qfullend,
% so the WHOLE question (stem, figures, options) is set across the page in reading order, then the two columns
% resume. The question's head - its number, stem and first figure (\qfullsplit marks the end of the head) -
% is kept on one page, so a number is never left at a page foot with its figure overleaf; the rest may break
% across pages normally, so no big gaps are left. Without a \qfullsplit, the whole question is kept together
% if it is under 45% of a page. Never write these yourself.
\newif\ifqf@full
\newif\ifqf@boxed
\newsavebox{\qf@box}
\newcommand{\qf@open}{\global\qf@fulltrue\par\global\qf@boxedtrue
  \setbox\qf@box\vbox\bgroup\hsize\textwidth\linewidth\textwidth}
\newcommand{\qf@emit}[1]{\par\egroup\global\qf@boxedfalse
  \ifdim\dimexpr\ht\qf@box+\dp\qf@box\relax<#1\textheight\noindent\usebox{\qf@box}\par\else\unvbox\qf@box\fi}
\newcommand{\qfullsplit}{\ifqf@boxed\qf@emit{0.85}\fi}
\newcommand{\qf@close}{\ifqf@boxed\qf@emit{0.45}\fi\global\qf@fullfalse}
\newcommand{\qfullbegin}{\end{multicols}\qf@open}
\newcommand{\qfullend}{\qf@close\begin{multicols}{2}}
\newcommand{\qfullstart}{\qf@open}
\newcommand{\qfulldone}{\qf@close}
% \qfig[scale]{id}: block figure, centred, directly where it is written (size set by the tool).
% Fallback if a wide figure is not inside a full-width question: break out of the columns in place - no floats.
\newcommand{\qfig}[2][1]{\qf@check{#2}%
  \ifqf@full\qf@blk{#1}{#2}\else
  \ifnum\@nameuse{qf@wide@#2}=1 \expandafter\@firstoftwo\else\expandafter\@secondoftwo\fi
  {\end{multicols}\begin{center}\qf@img{#1}{\textwidth}{#2}\end{center}\begin{multicols}{2}}%
  {\qf@blk{#1}{#2}}\fi}
\newcommand{\qf@blk}[2]{\par\smallskip{\centering\qf@img{#1}{\linewidth}{#2}\par}\smallskip}
% \qfigi[scale]{id}: inline figure (inside option lists / tables / minipages); never breaks columns.
\newcommand{\qfigi}[2][1]{\qf@check{#2}\qf@img[valign=c]{#1}{\linewidth}{#2}}
\makeatother
% structure macros - lint reads these to check completeness and figure placement
\newcommand{\papersec}[1]{\par\bigskip{\centering\Large\bfseries PAPER-#1\par}\medskip}
\newcommand{\partsec}[2]{\par\medskip\noindent{\large\bfseries PART-#1: #2}\par\smallskip}
\newcommand{\sect}[1]{\par\smallskip\noindent\textbf{#1}\par}
\newcommand{\qn}[1]{\par\smallskip\noindent\textbf{#1.}~}
\newcommand{\paraq}[2]{\par\smallskip\noindent\textbf{Paragraph for Questions #1 and #2}\par}
\newenvironment{qoptions}{\begin{itemize}[leftmargin=1.8em,itemsep=1pt,topsep=2pt,label={}]}{\end{itemize}}
\newcommand{\opt}[2]{\item[(#1)] #2}
% figure registry (generated by score2_tools.py finalize)
\qfigdef{p1_phy_s1_q02_1}{Internal Test - Test 02 (07 Jul 2021)-img/p1_phy_s1_q02_1.png}{96.0}{0}
\qfigdef{p1_phy_s1_q06_1}{Internal Test - Test 02 (07 Jul 2021)-img/p1_phy_s1_q06_1.png}{103.6}{0}
\qfigdef{p1_phy_s2_q07_1}{Internal Test - Test 02 (07 Jul 2021)-img/p1_phy_s2_q07_1.png}{154.5}{0}
\qfigdef{p1_phy_s2_q11_1}{Internal Test - Test 02 (07 Jul 2021)-img/p1_phy_s2_q11_1.png}{194.8}{0}
\qfigdef{p1_phy_s2_q12_1}{Internal Test - Test 02 (07 Jul 2021)-img/p1_phy_s2_q12_1.png}{66.5}{0}
\qfigdef{p1_phy_s3_q17_1}{Internal Test - Test 02 (07 Jul 2021)-img/p1_phy_s3_q17_1.png}{187.2}{0}
\qfigdef{p1_phy_s3_q18_1}{Internal Test - Test 02 (07 Jul 2021)-img/p1_phy_s3_q18_1.png}{245.2}{0}
\qfigdef{p1_chem_s1_q01_1}{Internal Test - Test 02 (07 Jul 2021)-img/p1_chem_s1_q01_1.png}{201.4}{0}
\qfigdef{p1_chem_s1_q02_A}{Internal Test - Test 02 (07 Jul 2021)-img/p1_chem_s1_q02_A.png}{54.7}{0}
\qfigdef{p1_chem_s1_q02_B}{Internal Test - Test 02 (07 Jul 2021)-img/p1_chem_s1_q02_B.png}{54.7}{0}
\qfigdef{p1_chem_s1_q02_C}{Internal Test - Test 02 (07 Jul 2021)-img/p1_chem_s1_q02_C.png}{68.9}{0}
\qfigdef{p1_chem_s1_q02_D}{Internal Test - Test 02 (07 Jul 2021)-img/p1_chem_s1_q02_D.png}{75.1}{0}
\qfigdef{p1_chem_s1_q03_1}{Internal Test - Test 02 (07 Jul 2021)-img/p1_chem_s1_q03_1.png}{125.5}{0}
\qfigdef{p1_chem_s1_q03_A}{Internal Test - Test 02 (07 Jul 2021)-img/p1_chem_s1_q03_A.png}{81.2}{0}
\qfigdef{p1_chem_s1_q03_B}{Internal Test - Test 02 (07 Jul 2021)-img/p1_chem_s1_q03_B.png}{81.2}{0}
\qfigdef{p1_chem_s1_q03_D}{Internal Test - Test 02 (07 Jul 2021)-img/p1_chem_s1_q03_D.png}{51.7}{0}
\qfigdef{p1_chem_s1_q04_1}{Internal Test - Test 02 (07 Jul 2021)-img/p1_chem_s1_q04_1.png}{128.9}{0}
\qfigdef{p1_chem_s1_q04_A}{Internal Test - Test 02 (07 Jul 2021)-img/p1_chem_s1_q04_A.png}{44.7}{0}
\qfigdef{p1_chem_s1_q04_B}{Internal Test - Test 02 (07 Jul 2021)-img/p1_chem_s1_q04_B.png}{44.1}{0}
\qfigdef{p1_chem_s1_q04_C}{Internal Test - Test 02 (07 Jul 2021)-img/p1_chem_s1_q04_C.png}{44.1}{0}
\qfigdef{p1_chem_s1_q04_D}{Internal Test - Test 02 (07 Jul 2021)-img/p1_chem_s1_q04_D.png}{44.1}{0}
\qfigdef{p1_chem_s1_q05_1}{Internal Test - Test 02 (07 Jul 2021)-img/p1_chem_s1_q05_1.png}{199.5}{0}
\qfigdef{p1_chem_s1_q05_A}{Internal Test - Test 02 (07 Jul 2021)-img/p1_chem_s1_q05_A.png}{41.0}{0}
\qfigdef{p1_chem_s1_q05_B}{Internal Test - Test 02 (07 Jul 2021)-img/p1_chem_s1_q05_B.png}{40.2}{0}
\qfigdef{p1_chem_s1_q05_C}{Internal Test - Test 02 (07 Jul 2021)-img/p1_chem_s1_q05_C.png}{35.1}{0}
\qfigdef{p1_chem_s1_q05_D}{Internal Test - Test 02 (07 Jul 2021)-img/p1_chem_s1_q05_D.png}{41.3}{0}
\qfigdef{p1_chem_s2_q07_A}{Internal Test - Test 02 (07 Jul 2021)-img/p1_chem_s2_q07_A.png}{106.5}{0}
\qfigdef{p1_chem_s2_q07_B}{Internal Test - Test 02 (07 Jul 2021)-img/p1_chem_s2_q07_B.png}{76.8}{0}
\qfigdef{p1_chem_s2_q07_C}{Internal Test - Test 02 (07 Jul 2021)-img/p1_chem_s2_q07_C.png}{57.6}{0}
\qfigdef{p1_chem_s2_q07_D}{Internal Test - Test 02 (07 Jul 2021)-img/p1_chem_s2_q07_D.png}{57.6}{0}
\qfigdef{p1_chem_s2_q08_A}{Internal Test - Test 02 (07 Jul 2021)-img/p1_chem_s2_q08_A.png}{95.1}{0}
\qfigdef{p1_chem_s2_q08_B}{Internal Test - Test 02 (07 Jul 2021)-img/p1_chem_s2_q08_B.png}{111.7}{0}
\qfigdef{p1_chem_s2_q08_C}{Internal Test - Test 02 (07 Jul 2021)-img/p1_chem_s2_q08_C.png}{139.0}{0}
\qfigdef{p1_chem_s2_q08_D}{Internal Test - Test 02 (07 Jul 2021)-img/p1_chem_s2_q08_D.png}{97.3}{0}
\qfigdef{p1_chem_s2_q09_A}{Internal Test - Test 02 (07 Jul 2021)-img/p1_chem_s2_q09_A.png}{31.5}{0}
\qfigdef{p1_chem_s2_q09_B}{Internal Test - Test 02 (07 Jul 2021)-img/p1_chem_s2_q09_B.png}{38.4}{0}
\qfigdef{p1_chem_s2_q09_C}{Internal Test - Test 02 (07 Jul 2021)-img/p1_chem_s2_q09_C.png}{31.2}{0}
\qfigdef{p1_chem_s2_q09_D}{Internal Test - Test 02 (07 Jul 2021)-img/p1_chem_s2_q09_D.png}{96.6}{0}
\qfigdef{p1_chem_s2_q10_A}{Internal Test - Test 02 (07 Jul 2021)-img/p1_chem_s2_q10_A.png}{132.9}{0}
\qfigdef{p1_chem_s2_q10_B}{Internal Test - Test 02 (07 Jul 2021)-img/p1_chem_s2_q10_B.png}{134.4}{0}
\qfigdef{p1_chem_s2_q10_C}{Internal Test - Test 02 (07 Jul 2021)-img/p1_chem_s2_q10_C.png}{123.5}{0}
\qfigdef{p1_chem_s2_q10_D}{Internal Test - Test 02 (07 Jul 2021)-img/p1_chem_s2_q10_D.png}{168.5}{0}
\qfigdef{p1_chem_s2_q11_1}{Internal Test - Test 02 (07 Jul 2021)-img/p1_chem_s2_q11_1.png}{197.3}{0}
\qfigdef{p1_chem_s2_q11_A}{Internal Test - Test 02 (07 Jul 2021)-img/p1_chem_s2_q11_A.png}{120.0}{0}
\qfigdef{p1_chem_s2_q11_3}{Internal Test - Test 02 (07 Jul 2021)-img/p1_chem_s2_q11_3.png}{79.6}{0}
\qfigdef{p1_chem_s2_q11_4}{Internal Test - Test 02 (07 Jul 2021)-img/p1_chem_s2_q11_4.png}{59.4}{0}
\qfigdef{p1_chem_s2_q11_D}{Internal Test - Test 02 (07 Jul 2021)-img/p1_chem_s2_q11_D.png}{120.6}{0}
\qfigdef{p1_chem_s3_q13_1}{Internal Test - Test 02 (07 Jul 2021)-img/p1_chem_s3_q13_1.png}{199.8}{0}
\qfigdef{p1_chem_s3_q14_1}{Internal Test - Test 02 (07 Jul 2021)-img/p1_chem_s3_q14_1.png}{214.9}{0}
\qfigdef{p1_chem_s3_q14_2}{Internal Test - Test 02 (07 Jul 2021)-img/p1_chem_s3_q14_2.png}{97.7}{0}
\qfigdef{p1_chem_s3_q14_3}{Internal Test - Test 02 (07 Jul 2021)-img/p1_chem_s3_q14_3.png}{58.5}{0}
\qfigdef{p1_chem_s3_q15_1}{Internal Test - Test 02 (07 Jul 2021)-img/p1_chem_s3_q15_1.png}{351.5}{1}
\qfigdef{p1_chem_s3_q16_1}{Internal Test - Test 02 (07 Jul 2021)-img/p1_chem_s3_q16_1.png}{245.2}{0}
\qfigdef{p1_chem_s3_q17_1}{Internal Test - Test 02 (07 Jul 2021)-img/p1_chem_s3_q17_1.png}{135.1}{0}
\qfigdef{p1_chem_s3_q18_1}{Internal Test - Test 02 (07 Jul 2021)-img/p1_chem_s3_q18_1.png}{58.5}{0}
\raggedcolumns
\makeatletter\@beginparpenalty=10000\makeatother
\newcolumntype{L}[1]{>{\raggedright\arraybackslash}p{#1}}
\begin{document}
\begin{center}{\Large\bfseries Internal Test --- Test 02 (07 Jul 2021)}\end{center}
\begin{multicols}{2}
\raggedcolumns
\papersec{1}
\partsec{1}{PHYSICS}
\sect{SECTION-I(i)}
\qn{1} A potentiometer wire has length $L$ and the emf of the driver cell is $E_0$ with negligible internal resistance. With a cell of emf $E$, the balance length is $\dfrac L4$. If the length of the potentiometer wire is doubled (without changing its diameter), the new balance length with the same cell will be
\begin{qoptions}
\opt{A}{$\dfrac L4$}
\opt{B}{$\dfrac L3$}
\opt{C}{$\dfrac L2$}
\opt{D}{$\dfrac{2L}{3}$}
\end{qoptions}
\qn{2} The potential difference through the $3\,\Omega$ resistor shown in figure :\par
\qfig{p1_phy_s1_q02_1}
\begin{qoptions}
\opt{A}{Zero}
\opt{B}{$1$ V}
\opt{C}{$3.5$ V}
\opt{D}{$7$ V}
\end{qoptions}
\qn{3} A cylindrical block of wood floats vertically with $80\%$ of its volume immersed in a liquid at $0^\circ$C. When the temperature of the liquid is raised to $62.5^\circ$C, the block just sinks in the liquid. The coefficient of cubical expansion of liquid is
\begin{qoptions}
\opt{A}{$1\times10^{-3}\,\mathrm{K^{-1}}$}
\opt{B}{$2\times10^{-3}\,\mathrm{K^{-1}}$}
\opt{C}{$3\times10^{-3}\,\mathrm{K^{-1}}$}
\opt{D}{$4\times10^{-3}\,\mathrm{K^{-1}}$}
\end{qoptions}
\qn{4} A metal ball immersed in alcohol weighs $W_1$ at $0^\circ$C and $W_2$ at $59^\circ$C. The coefficient of cubical expansion of metal is equal than that of alcohol. If the density of the metal is large compared to that of alcohol, then
\begin{qoptions}
\opt{A}{$W_1>W_2$}
\opt{B}{$W_1=W_2$}
\opt{C}{$W_1<W_2$}
\opt{D}{$W_2=\dfrac{W_1}{2}$}
\end{qoptions}
\qn{5} Two objects of masses $m$ and $4m$ are at rest at infinite separation. They move towards each other under mutual gravitational attraction. Then, at a separation $r$, which of the following is true?
\begin{qoptions}
\opt{A}{The total kinetic energy of the system is not zero.}
\opt{B}{The force between them is not zero.}
\opt{C}{The centre of mass of the system is at rest.}
\opt{D}{All the above are true.}
\end{qoptions}
\qn{6} A point charge $+q$ is projected from point A towards an infinitely long line charge having linear charge density $\lambda$ with kinetic energy $K_0$. The distance of the closest approach will be :\par
\qfig{p1_phy_s1_q06_1}
\begin{qoptions}
\opt{A}{$a\,e^{-\left(\frac{2\pi\varepsilon_0K_0}{q\lambda}\right)}$}
\opt{B}{$a\,e^{-\left(\frac{\pi\varepsilon_0K_0}{q\lambda}\right)}$}
\opt{C}{$a\left(1-\dfrac{2\pi\varepsilon_0K_0}{q\lambda}\right)$}
\opt{D}{$\left(\dfrac{2\pi\varepsilon_0K_0}{q\lambda}\right)$}
\end{qoptions}
\sect{SECTION-I(ii)}
\qn{7} In the circuit shown in figure, cells $E_1$ and $E_2$ have emfs $4$ V and $8$ V and internal resistances $0.5\,\Omega$ and $1\Omega$ respectively.\par
\qfig{p1_phy_s2_q07_1}
\begin{qoptions}
\opt{A}{The potential difference across $E_1$ is $4.25$ V}
\opt{B}{The potential difference across $E_1$ is $3.75$ V}
\opt{C}{The potential difference across $E_2$ is $8.55$ V}
\opt{D}{The potential difference across $E_2$ is $7.5$ V}
\end{qoptions}
\qn{8} If three bulbs of $40$ W, $60$ W and $100$ W are connected in series with a $200$ V power supply, then
\begin{qoptions}
\opt{A}{The potential difference will be maximum across the $40$ W bulb}
\opt{B}{$100$ W bulb will glow brightest}
\opt{C}{The $40$ W bulb has the maximum resistance}
\opt{D}{The current in the circuit is $1.2$ A}
\end{qoptions}
\qn{9} A satellite of mass $m$ is moving in a circular orbit of radius $r$ around a planet of mass $M$
\begin{qoptions}
\opt{A}{The magnitude of angular momentum with respect to the centre of the orbit is $m\sqrt{GMr}$, where $G$ is the gravitational constant.}
\opt{B}{The magnitude of the angular momentum is $mR\sqrt{2gr}$ where $g$ is the acceleration due to gravity on the surface of the planet.}
\opt{C}{The direction of angular momentum is parallel to the plane of the orbit.}
\opt{D}{The direction of angular momentum is perpendicular to the plane of the orbit.}
\end{qoptions}
\qn{10} A positive charge $q$ is fixed at the origin. An electric dipole with dipole moment $\vec p$ is placed along the x-axis far away from the origin with $\vec p$ pointing along the positive x-axis. The kinetic energy when it reaches a distance $x$ from the origin is $K$ and the magnitude of the force experienced by charge $q$ at this moment is $F$. Then
\begin{qoptions}
\opt{A}{$K$ veries as $\dfrac1x$}
\opt{B}{$K$ veries as $\dfrac1{x^2}$}
\opt{C}{$F$ veries as $\dfrac1{x^2}$}
\opt{D}{$F$ veries as $\dfrac1{x^3}$}
\end{qoptions}
\qn{11} An ellipsoidal cavity is made within a perfect conductor. A positive charge $q$ is placed in the cavity and A and B are two points on the surface of the cavity as shown in figure. Then\par
\qfig{p1_phy_s2_q11_1}
\begin{qoptions}
\opt{A}{Electric field near A in the cavity $=$ electric field near B in the cavity.}
\opt{B}{Charge density at A $<$ charge density at B}
\opt{C}{Potential at A $=$ potential at B}
\opt{D}{Total electric flux through the gaussiau surface shown is $\dfrac{q}{\epsilon_0}$}
\end{qoptions}
\qn{12} An uncharged conducting thin square plate of side $1$ m is placed in an uniform electric field of magnitude $200$N/C as shown in the figure. The field direction is perpendicular to the plane of plate. For this situation, mark out the correct statement (s). [In the options only big faces have been considered]:\par
\qfig{p1_phy_s2_q12_1}
\begin{qoptions}
\opt{A}{The net charge on each face of the plate is zero.}
\opt{B}{The net charge on each face of the plate is non-zero and having the magnitude equal to $1.77\times10^{-9}$C.}
\opt{C}{The net electric field intensity inside the plate is zero.}
\opt{D}{The electric field intensity at the surface of plate is discontinuous while potential is continuous.}
\end{qoptions}
\sect{SECTION-II}
\qn{13} A body is dropped from a height $h$ equal to $\dfrac R2$, where $R$ is the radius of the earth. If it will hit the surface of the earth with a speed $v=av_e$, where $v_e$ is the escape velocity from the surface of the earth. Find $a$
\qn{14} Assuming the earth to be a sphere of uniform mass density, find the percentage decrease in the weight of a body when it is taken to the end of a tunnel $32$ km below the surface of the earth. Radius of earth $=6400$ km
\qn{15} A neutral hydrogen molecule has two protons and two electrons. If one of the electrons is removed we get a hydrogen molecular ion $\left(\mathrm{H_2^+}\right)$. In the ground state of $\mathrm{H_2^+}$ the two protons are separated by roughly $1.5\,\mathring{\mathrm A}$ and the electron is roughly $1\,\mathring{\mathrm A}$ from each proton. What is the potential energy of the system? (in eV)
\qn{16} Two isolated metal spheres of radii $R$ and $2R$ are charged such that both have the same surface charge density $\sigma$. The spheres are located far away from each other. When they are connected by thin conducting wire, the new surface charge density on the bigger sphere will be $a\sigma$ then find $a$.
\qn{17} In the potentiometer circuit shows in figure find the value of $\ell$ when the galvanometer shows no deflection. The length of wire AB is $100$ cm and its resistance is $3\Omega$ (in cm)\par
\qfig{p1_phy_s3_q17_1}
\qn{18} In the circuit shown figure, the current through $4\Omega$ resistor is (in amp)\par
\qfig{p1_phy_s3_q18_1}
\partsec{2}{CHEMISTRY}
\sect{SECTION-I(i)}
\qn{1} Arrange the following acids in decreasing order of acidity :\par
\qfig{p1_chem_s1_q01_1}
\begin{qoptions}
\opt{A}{$X>Z>Y$}
\opt{B}{$X>Y>Z$}
\opt{C}{$Z>X>Y$}
\opt{D}{$Z>Y>X$}
\end{qoptions}
\qn{2} Which of the following compounds is most stable ?
\begin{qoptions}
\opt{A}{\qfigi{p1_chem_s1_q02_A}}
\opt{B}{\qfigi{p1_chem_s1_q02_B}}
\opt{C}{\qfigi{p1_chem_s1_q02_C}}
\opt{D}{\qfigi{p1_chem_s1_q02_D}}
\end{qoptions}
\qn{3} \qfig{p1_chem_s1_q03_1}
\begin{qoptions}
\opt{A}{\qfigi{p1_chem_s1_q03_A}}
\opt{B}{\qfigi{p1_chem_s1_q03_B}}
\opt{C}{$\mathrm{H{-}C{\equiv}C{-}H}$}
\opt{D}{\qfigi{p1_chem_s1_q03_D}}
\end{qoptions}
\qn{4} \qfig{p1_chem_s1_q04_1}
\begin{qoptions}
\opt{A}{\qfigi{p1_chem_s1_q04_A}}
\opt{B}{\qfigi{p1_chem_s1_q04_B}}
\opt{C}{\qfigi{p1_chem_s1_q04_C}}
\opt{D}{\qfigi{p1_chem_s1_q04_D}}
\end{qoptions}
\qn{5} For the reaction\par
\qfig{p1_chem_s1_q05_1}
\begin{qoptions}
\opt{A}{\qfigi{p1_chem_s1_q05_A}}
\opt{B}{\qfigi{p1_chem_s1_q05_B}}
\opt{C}{\qfigi{p1_chem_s1_q05_C}}
\opt{D}{\qfigi{p1_chem_s1_q05_D}}
\end{qoptions}
\qn{6} The IUPAC name of $\mathrm{CH_3{-}CH(CH_3){-}}$ $\mathrm{CH(CH_3){-}CH_2NH_2}$ is :
\begin{qoptions}
\opt{A}{2, 3 - Dimethylbutan-4-amine}
\opt{B}{2, 4--Dimethylbutan-1-amine}
\opt{C}{2,4-Dimethylbutan-4-amine}
\opt{D}{2, 3-Dimethylbutan-1-amine}
\end{qoptions}
\sect{SECTION-I(ii)}
\qn{7} Tautomerism is exhibited by :
\begin{qoptions}
\opt{A}{\qfigi{p1_chem_s2_q07_A}}
\opt{B}{\qfigi{p1_chem_s2_q07_B}}
\opt{C}{\qfigi{p1_chem_s2_q07_C}}
\opt{D}{\qfigi{p1_chem_s2_q07_D}}
\end{qoptions}
\qn{8} Which of the following reactions will give least substituted alkene ?
\begin{qoptions}
\opt{A}{\qfigi{p1_chem_s2_q08_A}}
\opt{B}{\qfigi{p1_chem_s2_q08_B}}
\opt{C}{\qfigi{p1_chem_s2_q08_C}}
\opt{D}{\qfigi{p1_chem_s2_q08_D}}
\end{qoptions}
\qn{9} Which of the following carbocations would rearrange ?
\begin{qoptions}
\opt{A}{\qfigi{p1_chem_s2_q09_A}}
\opt{B}{\qfigi{p1_chem_s2_q09_B}}
\opt{C}{\qfigi{p1_chem_s2_q09_C}}
\opt{D}{\qfigi{p1_chem_s2_q09_D}}
\end{qoptions}
\qn{10} Which of the following reactions are correct ?
\begin{qoptions}
\opt{A}{\qfigi{p1_chem_s2_q10_A}}
\opt{B}{\qfigi{p1_chem_s2_q10_B}}
\opt{C}{\qfigi{p1_chem_s2_q10_C}}
\opt{D}{\qfigi{p1_chem_s2_q10_D}}
\end{qoptions}
\qn{11} \qfig{p1_chem_s2_q11_1}
Select correct options ?
\begin{qoptions}
\opt{A}{\qfigi{p1_chem_s2_q11_A}}
\opt{B}{Y is \qfigi{p1_chem_s2_q11_3}}
\opt{C}{X is \qfigi{p1_chem_s2_q11_4}}
\opt{D}{\qfigi{p1_chem_s2_q11_D}}
\end{qoptions}
\qn{12} $y\ \xleftarrow[\text{Catalyst}]{\text{Lindlar's}}\ \mathrm{R{-}C{\equiv}C{-}R}\ \xrightarrow{\mathrm{Na/NH_3}}\ x$\par
$x$ and $y$ are geometrical isomers $\left(\mathrm{R{-}CH{=}CH{-}R}\right)$ of which type ?
\begin{qoptions}
\opt{A}{$x$ is trans}
\opt{B}{$y$ is cis}
\opt{C}{$x$ and $y$ both are trans}
\opt{D}{$x$ is cis, $y$ is trans}
\end{qoptions}
\sect{SECTION-II}
\qn{13} Find out number of aromatic compounds or ions from the following\par
\qfig{p1_chem_s3_q13_1}
\qn{14} How many of the following do not have plane of symmetry ?\par
\qfig{p1_chem_s3_q14_1}
\qfig{p1_chem_s3_q14_2}
\qfig{p1_chem_s3_q14_3}
\qfullbegin
\qn{15} How many of the following addition reactions are syn addition\par
\qfig{p1_chem_s3_q15_1}
\qfullend
\qn{16} Examine the structural formulas of following compounds and find out how many compounds will produce $\mathrm{CO_2}$ on oxidative ozonolysis.\par
\qfig{p1_chem_s3_q16_1}
\qn{17} What is the degree of unsaturation in\par
\qfig{p1_chem_s3_q17_1}
\qn{18} Find out number of benzylic hydrogen in\par
\qfig{p1_chem_s3_q18_1}
\partsec{3}{MATHEMATICS}
\sect{SECTION-I(i)}
\qn{1} Let $f:(1,\infty)\to(0,\infty)$ be continuous strictly decreasing function with $\displaystyle\lim_{x\to\infty}\dfrac{f(4x)}{f(8x)}=1$. Then $\displaystyle\lim_{x\to\infty}\dfrac{f(6x)}{f(8x)}=$
\begin{qoptions}
\opt{A}{$\dfrac12$}
\opt{B}{$\dfrac23$}
\opt{C}{$\dfrac34$}
\opt{D}{$1$}
\end{qoptions}
\qn{2} The value of $a$ for which $ax^2+\sin^{-1}\left(x^2-2x+2\right)+\cos^{-1}\left(x^2-2x+2\right)=0$ has a real solution,
\begin{qoptions}
\opt{A}{$\dfrac\pi2$}
\opt{B}{$\dfrac{-\pi}2$}
\opt{C}{$\dfrac2\pi$}
\opt{D}{$-\dfrac2\pi$}
\end{qoptions}
\qn{3} The equation $2\cos^{-1}x+\sin^{-1}x=\dfrac{11\pi}6$ has
\begin{qoptions}
\opt{A}{No solution}
\opt{B}{Only one solution}
\opt{C}{Two solution}
\opt{D}{Three solution}
\end{qoptions}
\qn{4} A continuous function $f(x)$ on $R\to R$ satisfies the relation $f(x)+f(2x+y)+5xy=f(3x-y)+2x^2+1\ \ \forall x,y\in R$, then $f(x)$ is
\begin{qoptions}
\opt{A}{$1-x^2$}
\opt{B}{$1-\dfrac{x^2}2$}
\opt{C}{$1-2x^2$}
\opt{D}{$2-x^2$}
\end{qoptions}
\qn{5} $f(x)=-1+|x-1|,-1\le x\le3$\par
$g(x)=2-|x+1|,-2\le x\le2$\par
then for $x\in(0,1)$, $fog(x)$ equals
\begin{qoptions}
\opt{A}{$x-1$}
\opt{B}{$1-x$}
\opt{C}{$x+1$}
\opt{D}{$-(x+1)$}
\end{qoptions}
\qn{6} If $\sqrt{x+y}+\sqrt{y-x}=5$, then $\dfrac{d^2y}{dx^2}=$
\begin{qoptions}
\opt{A}{$\dfrac25$}
\opt{B}{$\dfrac4{25}$}
\opt{C}{$\dfrac2{25}$}
\opt{D}{$\dfrac1{25}$}
\end{qoptions}
\sect{SECTION-I(ii)}
\qn{7} $\displaystyle\lim_{n\to\infty}\left\{\left(\sqrt5+2\right)^n\right\}\cdot\cos\left(\pi\left[\left(\sqrt5+2\right)^n\right]\right)$ is equal to\par
(where $n\in I^+$ and $[\bullet],\{\bullet\}$ denotes greatest integer function and fractional part function respectively)
\begin{qoptions}
\opt{A}{$0$ when $n$ is even}
\opt{B}{$-1$ when $n$ is even}
\opt{C}{$-1$ when $n$ is odd}
\opt{D}{$0$ when $n$ is odd}
\end{qoptions}
\qn{8} $f$ is continuous function in $[\,a,b]$, $g$ is a continuous function in $[b,c]$. A function $h(x)$ is defined as\par
$h(x)=\begin{cases}f(x)&\text{for }x\in[a,b)\\ g(x)&\text{for }x\in(b,c]\end{cases}$, If $f(b)=g(b)$ , then
\begin{qoptions}
\opt{A}{$h(x)$ has a removable discontinuity at $x=b$}
\opt{B}{$h(x)$ may or may not be continuous in $[a,c]$}
\opt{C}{$h\left(b^-\right)=g\left(b^-\right)$ and $h\left(b^+\right)=f\left(b^-\right)$}
\opt{D}{$h\left(b^+\right)=g\left(b^-\right)$ and $h\left(b^-\right)=f\left(b^+\right)$}
\end{qoptions}
\qn{9} If $2^{\left(\cos^{-1}x\right)^2}\sqrt{\left(\sin^{-1}y\right)^4-2\left(\sin^{-1}y\right)^2+2}=1$, then the value of $(x-y)$ can be :
\begin{qoptions}
\opt{A}{$1-\dfrac\pi2$}
\opt{B}{$1+\dfrac\pi2$}
\opt{C}{$1-\sin1$}
\opt{D}{$1+\sin1$}
\end{qoptions}
\qn{10} Which of the following is periodic ?
\begin{qoptions}
\opt{A}{$f(x)=\{x\}$,where $\{x\}$ denotes the fractional part of the number $x$}
\opt{B}{$f(x)=\cot(x+7)$}
\opt{C}{$f(x)=1-\dfrac{\sin^2x}{1+\cot x}-\dfrac{\cos^2x}{1+\tan x}$}
\opt{D}{$f(x)=x+\sin x$}
\end{qoptions}
\qn{11} If `a' is a real number for which $f(x)=\ln\left(3\cos^{-1}\left(\dfrac{3x}7\right)-\pi\right)$ is defined then the possible value(s) of $[a]$ is (where $[.]$ denotes greatest integer function)
\begin{qoptions}
\opt{A}{$-2$}
\opt{B}{$-1$}
\opt{C}{$1$}
\opt{D}{$0$}
\end{qoptions}
\qn{12} Let $f(x)=\begin{cases}\dfrac{\left(2^x-1\right)^2}{\tan x\cdot\ln(1+x)}&;\ x\ne0\\[2ex] \ln4&;\ x=0\end{cases}$\par
then $f(x)$ is
\begin{qoptions}
\opt{A}{Continuous at $x=0$}
\opt{B}{Non-removable discontinuity at $x=0$}
\opt{C}{$f(x)$ has a removable discontinuity at $x=0$}
\opt{D}{$\displaystyle\lim_{x\to0}f(x)$ exists}
\end{qoptions}
\sect{SECTION-II}
\qn{13} Define functions $f_0,f_1,\ldots\ldots\ldots,f_n,\ldots.$ to satisfy the following relations : $f_0=2,\ \dfrac{d}{dx}\left(f_n\right)=f_{n-1}$.\par
If for $n\ge1$, $f_n(0)=0$, compute $\displaystyle\sum_{n=0}^{\infty}f_n\left(\ell n(2019)\right)$
\qn{14} Let $f(x)\ \&\ g(x)$ be continuous function for all $x\in R$. If $\displaystyle\lim_{x\to3}\dfrac{f(x)}{x-3}=\lim_{x\to f(3)}\dfrac{g(x)}{x^2}=5$ then $\displaystyle\lim_{x\to g(f(3))}\dfrac{\tan10x}{x}=$
\qn{15} Let $f(x)=\dfrac{e^x-e^{-x}}2$ and if $f(g(x))=x$ then $g\left(\dfrac{e^{1002}-1}{2e^{501}}\right)$ is equal to
\qn{16} If $2<x^2<3$ then the number of positive roots of $\left\{\dfrac1x\right\}=\left\{x^2\right\}$, denotes the fractional part of $x$, is
\qn{17} The number of solutions of the equation $\sin\pi x=\left|\log_e|x|\right|$ is
\qn{18} Let $f(x)=3x^{10}-7x^8+5x^6-21x^3+3x^2-7$\par
$265\left(\displaystyle\lim_{h\to0}\dfrac{h^4+3h^2}{\left(f(1-h)-f(1)\right)\sin5h}\right)=$
\end{multicols}
\end{document}
